\vglue -2cm
\renewcommand{\arraystretch}{1.28}
\noindent
\hspace{-0.5cm}\begin{tabularx}{1.059\linewidth}{|c|Y|c|c|c|}\hline
 \multirow{4}{*}{\parbox[b]{3cm}{\vspace{0.1cm}\centering\includegraphics[scale=0.3]{udea}\\
\scriptsize{{\fontfamily{ptm}\selectfont Universidad de Antioquia}}\\
\tiny{{\fontfamily{ptm}\selectfont Vicerrectoría de Investigación}}
}
}  & \textbf{\large VICERRECTORÍA DE INVESTIGACIÓN} & \multicolumn{3}{c|}{{\large Actualización}}\\
\cline{2-5}
 &\textbf{FICHA TÉCNICA}  & {\large Día} & {\large Mes} & {\large Año}\\
\cline{3-5}
 &  &  &  & \\[-13pt]
 & \textbf{\small PROYECTOS DE INVESTIGACIÓN} &  &  & \\[-5pt]
 & \textbf{2009} &  &  & \\
\hline
\end{tabularx}
%
\renewcommand{\arraystretch}{1}

\bigskip{}
\noindent{}
La ficha técnica es el resumen del proyecto, con una extensión máxima de seis páginas sin contar anexos. El investigador debe tener especial cuidado para que la ficha técnica contenga la información básica necesaria acerca del proyecto, puesto que el jurado final solo recibe la mencionada ficha. El proyecto de investigación quedará en archivo en el centro de investigaciones.



\begin{center}
  \textbf{\large INFORMACIÓN GENERAL}
\end{center}

\renewcommand{\arraystretch}{1.71}
\noindent
\hspace{-0.5cm}\begin{tabularx}{1.059\linewidth}{|l|l|}\hline
\multicolumn{2}{|Y|}{\cellcolor[gray]{.8}\textbf{Implicaciones de una simetría $U(1)_H$ en experimentos de rayos cósmicos y de aceleradores}}\\\hline
\multicolumn{2}{|Y|}{}\\\hline
\cellcolor[gray]{.8}\textbf{Facultad, Escuela o Instituto}&Facultad de Ciencias Exactas y Naturales\\\hline
\cellcolor[gray]{.8}\textbf{Grupo que avala el proyecto}& Fenomenología de Interacciones Fundamentales (GFIF)\\\hline
\cellcolor[gray]{.8}\textbf{Código Colciencias del Grupo si aplica}&COL0008423\\\hline
\cellcolor[gray]{.8}\textbf{Nombre del Investigador Principal}&Diego A. Restrepo Qunitero\\\hline
\cellcolor[gray]{.8}\textbf{Nombre del Coinvestigador}&William Ponce\\\hline
\cellcolor[gray]{.8}\textbf{Documento de Identidad N${}^{\text{\textbf{o}}}$}&98554575\\\hline
\cellcolor[gray]{.8}\textbf{Convocatoria}&Sostenibilidad 2009-2010 de GFIF\\\hline
\cellcolor[gray]{.8}\textbf{Duración del proyecto (en meses)}&18 meses\\\hline
\end{tabularx}

\vspace{-0.02cm}
\renewcommand{\arraystretch}{1.038}
\noindent
\hspace{-0.5cm}\begin{tabularx}{1.059\linewidth}{|l|r|}
\multicolumn{2}{|Y|}{\cellcolor[gray]{.8}\textbf{DATOS DEL PRESUPUESTO}}\\[7pt]\hline
\textbf{Recursos propios del Grupo}\hspace{6.2cm}\quad&{\$ 14,000,000,\footnotesize{00}}\\\hline
\textbf{Universidad} \footnotesize{(recursos en especie y recursos  frescos)}&{\$ 43,750,000\footnotesize{00}}\\\hline
\textbf{Otras entidades} \footnotesize{(si aplica):}&{\$ 0,\footnotesize{00}}\\\hline
\textbf{Valor total del proyecto:}&{\$ 57,750,000,\footnotesize{00}}\\\hline
\end{tabularx}
\renewcommand{\arraystretch}{1}

\begin{center}
  \textbf{\large DESCRIPCIÓN DEL PROYECTO}
\end{center}

%%% Local Variables: 
%%% mode: latex
%%% TeX-master: t
%%% End: 
